\chapter{Introduction}
\label{chp:intro}

% Answer the following question

% What is the context of your thesis?  What area of computer science is it located in?
% What is the problem that you are addressing?
% How do you try to solve the problem? -> implement a Rust adapter for ECCO
% How do you evaluate your solution? -> 
% What is the result of your evaluation?
% How is the thesis structured in the following?
Modern software systems are rarely developed as a single, fixed program.  
Instead, many systems are configurable, meaning that the same code base can be
compiled into different software variants depending on hardware, platform, or
selected features.
A common way to implement such configurability is through conditional code
guarded by presence conditions, for example using preprocessor directives such
as \texttt{\#if} and \texttt{\#ifdef}.  
Presence conditions determine in which configurations a particular code fragment
is included and therefore directly influences the behavior of the resulting
software variants.

This bachelor thesis is situated in the area of software engineering, with a
focus on configurable software systems and variability analysis.
It studies how changes to conditional code, such as presence conditions
e.g. \texttt{\#ifdef}, affect software systems over time by analysing large
software repositories.
In standard version control systems, changes are represented as textual diffs
that show added or removed lines of code.  
While this representation is sufficient for many development tasks, it is often
insufficient for configurable systems.  
In such systems, even a small textual change can have a large semantic impact if
it modifies a presence condition.

For example, a minor edit to a conditional expression may silently expand or
restrict the set of configurations in which a code fragment is active.  
These effects are not visible in standard diffs and are often undocumented in
commit messages.  
As a result, developers and tools may overlook unintended changes to variability,
which can lead to configuration of specific faults that are difficult to detect.
The problem addressed in this thesis is therefore the limited ability of
traditional diffs to explain the impact of changes to presence conditions.  
In particular, the thesis investigates how developer mistakes related to
variability can be identified by analyzing commits at the level of variability
semantics rather than raw text differences.

To address this problem, this thesis builds on the DiffDetective tool to explicitly
identify changes in variability. 
DiffDetective analyzes software evolution by mapping source code and annotations
to structured representations and by classifying edits according to their effect
on variability.  
Instead of treating all changes as simple additions or deletions, DiffDetective
distinguishes different kinds of variability edits, such as refactoring and
reconfiguration.
Based on this foundation, this thesis contributes a custom analysis tool called
\texttt{MyAnalysis}.  
The purpose of \texttt{MyAnalysis} is to identify and store variability edits that
are relevant for understanding developer mistakes.  
The recorded results are later analyzed using a topic-based approach, which
groups refactoring and reconfiguration edits and collecting their diff messages.
This makes it easier to study variability related diffs in more detail and to
reason about recurring patterns in how such changes occur.
The analysis focuses mainly on two classes of edits: refactoring and
reconfiguration.  
Refactoring describes changes where the structure of a presence condition is
modified while its meaning remains the same.  
Reconfiguration describes changes where the artifact remains unchanged, but the
set of configurations in which it is included changes in a non-trivial way.

In addition, \texttt{MyAnalysis} extracts commit-level information such as commit
identifiers and commit messages.  
A simple keyword-based heuristic is used to flag commits that are likely related
to bug fixes.  
This allows the analysis to connect variability edits with developer intent as
expressed in commit messages.

The solution is evaluated by applying \texttt{MyAnalysis} to a real-world
configurable software system.  By analysing how often variability related diffs occur
in practice and by examining which types of diffs are most common and studying the distribution of different variability edit types across large
software repositories, the evaluation makes it possible to assess the overall
extent of the problem.  
This approach allows us to determine whether changes to presence conditions occur
frequently and whether they represent a more substantial issue than initially
expected.
The results show that the analysis is able to distinguish between different types
of variability changes and provides a more transparent view of how presence
conditions evolve over time.  
By going beyond traditional textual diffs, the approach makes variability changes
visible that would otherwise be difficult to recognise in standard version
control output.

In the following chapters, this thesis is structured as follows.  
The next chapter provides background on configurable software and presence
Conditions and explains the motivation for using variability aware
differencing.  
The subsequent chapter describes the design and implementation of
\texttt{MyAnalysis}.  
After that, the evaluation chapter presents the study and discusses the results.
Finally, the thesis concludes with a summary of the main findings and an outline
of possible directions for future work.
% the following is to see how the individual pages look like (we assume a 2-page book so they should look differently)
\clearpage

trying with novel page

\clearpage
and another one